\subsection{次の接着計算に使う二つの局所障害}
\label{ng16:local-obstructions}

二価 core $i$ の二つの高近傍を $j,k$ とする．$q_{jk}=1$ なら，
core $i$ の二高分母標準形は必ず正の枝である．従って
\[
 S>r,\qquad rT>(r+1)S,\qquad L=T-S+1\geq3.
\]
この結論に対向 $D$ の相殺は不要であり，core $j,k$ の高分母次数を制限しない．
負の枝 $S=T=w<r$ を用いるときは，その core を含む高分母三角形が必要である．
一般の三高分母中心に二高分母の真約数補題を適用する代わりに，
中心側では signed divisors と高辺の合同式を直接比較する．
負の枝の $H=r-w$，$c=s-x$，$e=z-y$ に対して
\[
 Ece=e-c+1,\qquad E=1-\frac{w^2}{H(H+1)},
 \qquad -1<E<\tfrac12,\quad E\ne0
\]
となり，中心間が低辺なら reverse 和が $E-1$ となることが主な障害である．
例外 $e=1$ も node 数 $2+rw/H>6$ により排除される．

全低分母 core $*$ がある場合には，残る四頂点を $I$ とし
\[
 x_{ij}=p^i_{j*},\qquad r_{ij}=p^i_{*j}
\]
と置く．四辺 $i*$ の $D$ 相殺，局所の行和，および高辺の整数 clutching から
\[
 \sum_{i\ne j\in I}x_{ij}=\sum_{i\ne j\in I}r_{ij}=-4,
 \qquad
 \sum_{\substack{i<j\in I\\q_{ij}=1}}
 (x_{ij}+x_{ji}-r_{ij}-r_{ji})=0
\]
を得る．全低 core の自由な entries を固定しなくても使える総和条件である．
diamond では和に低辺が一つだけ残り，二価 core の低低 entries が
反対向き・同じ $L\in\{1,3,4\}$ にそろう．三つの場合は全て排除済みである．
$K_4$ 型ではこの和は空になるが，四本の $D$ 相殺を頂点別に保持すると，
後述する通常辺上の流れの式を得る．


\subsection{共通低 puncture における通常接着の上界}
\label{ng16:ordinary-bound}

$q_{ij}>1$，$q_{ik}=q_{jk}=1$ とし，高辺 $ij$ の puncture $k$ が通常の反対符号接着なら
\[
 p^i_{jk}=-p^j_{ik},\qquad p^i_{kj}=-p^j_{ki}=u,\qquad1\le |u|\le4.
\]
core $k$ の低低 entries を $a=p^k_{ij},b=p^k_{ji}$ とすれば $ab=0$ であり，
二つの node 数 $n,m\in\{0,\ldots,5\}$ に対して
\[
 a=1-n-u,\qquad b=1-m+u,\qquad
 0=-\{u-(1-n)\}\{u-(m-1)\}
\]
となるためである．高分母次数と対向 $D$ の相殺には依存しない．
node 総数が同じ $N\ge2$ の模型では，同じ証明で $|u|\le N-1$ を得る．
最大・単位約数の同符号例外をこの上界に含めてはいけない．

\subsection{一高分母の自由度を保持する標準形}
\label{ng16:leaf-data}

一高分母 core は，高 ray $A$ と低 rays $B,C,D$ の適切な標識により
\[
\begin{gathered}
 m=th(h+1),\quad q=m+1,\quad u=th,\quad v=t(h+1),
 \quad h,t\geq1,\\
 P=\begin{pmatrix}
 0&-m&-h-1&h\\
 -1&0&0&0\\
 -u&u-1-c&0&c\\
 v&-v-1-d&d&0
 \end{pmatrix},\qquad c,d\in\ZZ,\quad cd=0
\end{gathered}
\]
と書け，全ての整数 choices に正の weights が存在する．係数は
\[
 (D_A,D_B,D_C,D_D)=(-m-t,\ m+t+1+c+d,\ h-d,\ -h-1-c).
\]
$q=3$ では二つの負の reverse entries が一致し得る．
高 row だけを固定して低低 entries を一つの模型に置き換えず，
この $c,d$ の自由度と必要な低 ray の交換を保持する．

二高分母の低係数は
\[
 D_C=y-x+T-S,\qquad D_D=s-z-1
\]
であり，同時に負とはならない．一高分母の負 reverse を持つ ray $E$ は $B$ または $C$ である．
別の低 ray $J$ から $p_{JE}\leq-1$ を要求される場合，
$E=C$ なら $D_E>0$，残る低 ray から $E$ への entry は零となる．
$E=B$ かつ $D_E\leq0$ なら，残る低 ray からの entry は
$t(h+1)^2-D_E\geq4$ である．
係数相殺と node 条件を組み合わせるときに使える二分である．

\subsection{二葉付き三角形の完全な有限化から再利用する式}
\label{ng16:two-leaf-certificate}

$H=\{01,02,12,03,14\}$ は全排除済みである．再利用するのは，
未知の中心分母を走査せず，葉の自由度を保存したまま有限化する次の手順である．
中央 core $2$ の交差標識 $(A,B,C,D)=(0,1,4,3)$ と正の枝に対して，
共通低 puncture の上界は
\[
 1\le r\le4,\qquad z=1
\]
を同時に要求する．$T\ge5$ により辺 $12$ は正の最大・単位約数の例外となり，
\[
 b=T+1,\quad p^1_{23}=T,\quad p^1_{32}=1,\quad
 p^3_{12}=0,\quad n_{12;3}=n_{13;2}=0
\]
を得る．$r>1$ なら，辺 $02$ には通常接着が強制される．

$T=(r+1)y$，$E=D^2_4$ と置くと，
\begin{align*}
 (r+2)sS-r(rs+1)y&=2r(s-1),\\
 F_r(s)S&=r(rs+1)E+2r(r+2)(s-1),\\
 F_r(s)&=-rs^2+(4r+3)s-r,\\
 E&=(r+2)y-\frac{s+r}{r}S.
\end{align*}
$r=1,\ldots,4$ では $F_r(s)\ne0$ であり，$E\ge-2$ なら $s\le9$ となる．
より一般に，整数の正枝で $z=1$ なら，任意の正整数 $r,s$ でも $F_r(s)\ne0$ である．
実際，$F_r(s)=0$ は $r=3s/(s^2-4s+1)$ を与え，正整数解は $(r,s)=(12,4)$ だけだが，
上の係数式は $E=-12/7$ を要求して整数性に反する．
従って $(r,s,E)$ は $S$，続いて $y,T$ を一意に復元する．
また $E\ge-M$，$M\in\ZZ_{\ge0}$ なら
\[
 E<\left(4+\frac4r-s\right)\frac Sr,\qquad S>r
\]
から $s\le M+7$ を得る．これは他の係数下界が得られた場合にも再利用できる．
$S\ge4$ の下で，通常の $C/D$ 表現を持つ葉 $4$ では，
$-2\le E\le7$ の枝と，$r=1$，$D^4_3=E+\kappa$，
$\kappa\in\{2,5,8,10\}$ の枝に分かれる．
後者でも葉 $3$ の一方の低標識なら $th\le4$，
他方なら $h\le12,t\le6$ が証明される．
小さい $S$ の共鳴族 $(r,S,T)=(2,3,6)$ では，$s$ は正の偶数であり，
$s\ge8$ は他葉に向かう二つの正の $D$ が相殺しないため排除される．

この有限化は三十一中央模型・六十六生成 witness を与え，
全葉を復元すると十中央模型・五十二低三角形だけが残る．
葉 $4$ の ray $2=B$ を含む別の有限化は，
三十三中央模型・六十二 witness から四中央模型・七組を与える．
この七組は五十二組の部分集合であり，個数を加えてはいけない．
分母三の同じ標識付き行列は，一度だけ数える．
$B$ 表現の例外を省くことも，低低自由度 $c,d$ を固定することも許されない．

\paragraph{既知三行だけによる中心の延長障害．}
三つの隣接 cores の行列が既知なら，中心 $i$ の row $j$ は
\[
 p^i_{jk}=q_{ij}(1-n_{ij;k})-p^j_{ik}
\]
で復元される．非負 node 数の三分割は二十一通りなので，既知三行は全 $21^3$ 通りである．
三つの既知 rays の pair equations は
\[
 p^i_{jk}p^i_{kj}=q_{ij}q_{ik}-1
\]
である．五十二組のうち五十一組は中心 $0$，残る一組は中心 $1$ において，
この三行条件だけで候補が零となる．
未知の中心間分母，第四行，中心での $D$ 相殺，正の weights の判定は，
この最後の矛盾に必要ない．
この手順は，他の残存 graph で複数の隣接行が既知となる場合にも使える．

\subsection{\texorpdfstring{$K_4$}{K4} の低 hub と通常辺上の流れ}
\label{ng16:k4-flow}

全低 core を $4$ とし，$I=\{0,1,2,3\}$ の高 $K_4$ で
\[
 R_{ij}=p^i_{4j},\quad X_{ij}=p^i_{j4},\quad V_{ij}=p^4_{ij},
 \quad N_{ij}=1-R_{ij}-V_{ij}
\]
と置く．各行の node 予算は $\sum_jN_{ij}=5$ である．
通常辺では $R_{ji}=-R_{ij}$，$1\le|R_{ij}|\le4$，
例外辺では $R$ の対が $\pm\{1,q-1\}$ となる．
hub に接する $D$ 相殺を頂点別に保持すると
\[
 \sum_{j\ne i}(X_{ij}-R_{ji})
 =\sum_{j\ne i}N_{ji}-5
\]
を得る．例外辺の寄与は零であり，通常辺の寄与は
\[
 \operatorname{sgn}(R_{ij})W_{ij},\qquad
 W_{ij}=\frac{q_{ij}-1}{|R_{ij}|}+|R_{ij}|
\]
である．proper cut の発散の絶対値は高々十なので，通常 graph の橋では
$W\le10$，従って $q\le25$ となる．

例外辺数は二・四・六のいずれかであるが，三場合は全て排除された．
二例外は，負の大きい reverse を持つ頂点から見た正例外の位置による三場合で，
hub の対向 $D$ 式と正 flow の下界が矛盾する．三高分母の一般分類は不要である．

四例外の再利用できる有限化は，例外の大きい reverse $m=q-1$ がその大きい端の
一つの行だけに現れることを用いる．異符号二項の差を固定した無限方向を
\[
 m_1=h+\max(\delta,0),\qquad m_2=h+\max(-\delta,0),\qquad h\in2\ZZ_{>0}
\]
と書き，$h=2,4$ と全 $h\ge6$ に分ける．最後の半直線では大きい端の
$V=1-R-N$ が必ず非零なので，pair 条件は単位端の node を正例外で零，負例外で二に固定する．
この後の row，node，hub pair，低辺の $D$ 条件は $h$ に依存しない．
従って $h=6$ は無限型全体の厳密な代表であり，分母cutoffではない．
通常二辺が端点非共有の場合の百十二型では候補零，共有の場合の二百四型では候補六となり，
残る無限型は零である．全六候補に
\[
 (q_1,q_2,q_3)=(5,3,3),\quad r=(-4,1,2),\quad x=(-1,2,1)
\]
のcoreが現れ，残る局所条件は $6A^2+17A+14=0$，判別式 $-47$ となる．
これは実数の高高entriesも許さない．

有限化と独立監査には，reverse-first と node-first の二方向がある．
後者は全 $21^4$ node 行列から通常forestのflowを決め，例外分母の行方程式を別に解く．
両方法は六候補の全行列で一致した．次のgraphでは，一方の表現で無限方向が見えにくい場合に
他方から復元し，同じ有限候補へ至るかを調べる価値がある．
全低coreを持つsupportはこれで全分母一に限られるが，一高分母hubへ移すと
pair productsとrow予算が変わるため，補正を導かずに同じ式を適用してはいけない．

全分母一のatlasや一高分母hubとの比較のため，全低 hub 自身の自由度の表示も残す．
$V_{ij}V_{ji}=0$，行和 $-1$ から，各行に負の非零 entry が少なくとも一つある．
非零 support は四，五，六辺であり，出次数列はそれぞれ
$(1,1,1,1)$，$(2,1,1,1)$，および $(3,1,1,1)$ または $(2,2,1,1)$ である．
出次数一の entry は $-1$，出次数二の行は $(t,-1-t)$ なので，
有限個の有向 support と高々二整数の族として扱える．
特定の全 $-1$ 模型に固定せず，正の weights が必要な段階で cycle 条件を課す．
\subsection{部分的な低辺条件だけでは完成しない二つの模型}
\label{ng16:partial-patches}

次の例は，部分的な必要条件の非空性を示す資料である．大域 atlas の存在例ではない．
$K_{2,3}$ の三つの二価 core を cyclic に標識し，canonical 高 ray $A$ が
一方の中心を向く場合を $+$，他方を向く場合を $-$ とする．
符号 $(+,+,-)$ に対し，canonical データを
\[
 (r,s,z,S,T)=(4,4,5,2,2),\quad(5,10,2,3,3),\quad(5,10,2,3,3)
\]
と取る．三低辺の node triples は $(0,3,2),(2,2,1),(0,2,3)$ だが，
同じ順の対向低係数和は $7,-10,5$ であり，相殺しない．
中心二 core も与えていない．本版の $K_{2,3}$ 排除は，中心間の低辺と高辺の clutching を加えて成立する．
具体的な全行列・係数・node の再計算は，保存した v14 の node 検算器に含まれる．

diamond の二価 cores $2,3$ と全低 core $4$ についても，
三つの相互低辺の node 条件と $D$ 相殺を全て満たす部分模型がある．
行列の添字順をそれぞれ $(0,1,3,4),(0,1,2,4),(0,1,2,3)$ として
\[
 P_2=\begin{pmatrix}0&-6&1&-2\\-8&0&-2&3\\6&-3&0&-4\\-3&2&0&0\end{pmatrix},\quad
 P_3=\begin{pmatrix}0&-14&-2&5\\-7&0&2&-4\\-5&4&0&0\\2&-2&-1&0\end{pmatrix},
\]
\[
 P_4=\begin{pmatrix}0&0&-1&0\\2&0&0&-3\\0&-1&0&0\\-3&0&2&0\end{pmatrix},
 \quad q_2=(7,7,1,1),\quad q_3=(11,9,1,1),\quad q_4=(1,1,1,1).
\]
係数は順に
\[
 (-2,0,0,2),\qquad(-1,3,0,-2),\qquad(0,0,-2,2)
\]
であり，低辺 $23,24,34$ の node triples はそれぞれ
$(0,0,5)$，$(4,0,1)$，$(2,3,0)$ である．
全てに正の weights がある．しかし全低 hub の残差
\[
 p^2_{34}+p^3_{24}-p^2_{43}-p^3_{42}=-3
\]
は零でないので，二つの中心 core を加えた大域完成は不可能である．
行和・pair products・正 weights・係数・nodes の厳密検算は
\nolinkurl{diamond_obstruction_verification_v15.json} に保存した．
部分模型の存在を保ちながら，追加する大域条件のどこで矛盾するかを区別するための例である．
